\section{From subgroup intervals to extensions}
\label{ext:section}

We now apply the decorated lattice of Proposition~\ref{lat:decorated-lattice}.
The reductions in this section use its tests in both orientations.  The
simple group attached to the top of the original group interval will be
denoted by $T$.  The simple group attached to an interval of extension
domains will later be denoted by $R$; these groups need not be isomorphic.

\begin{lemma}\label{ext:almost-simple}
Let $T$ be a finite nonabelian simple group and let
$\operatorname{Inn}(T)\le H\le\operatorname{Aut}(T)$.
Every nontrivial normal subgroup of $H$ contains $\operatorname{Inn}(T)$.
In particular, $H$ has trivial soluble radical and has characteristic
socle $\operatorname{Inn}(T)$.
\end{lemma}
\begin{proof}
An automorphism $\sigma$ centralizing every inner automorphism satisfies
$c_{\sigma(t)}=c_t$ for every $t\in T$.  Since $Z(T)=1$, this implies
$\sigma(t)=t$ for every $t$, so
$C_{\operatorname{Aut}(T)}(\operatorname{Inn}(T))=1$.
If $K\trianglelefteq H$, then
\[
 [K,\operatorname{Inn}(T)]\le K\cap\operatorname{Inn}(T).
\]
Simplicity makes the intersection either trivial or all of
$\operatorname{Inn}(T)$.  In the first case the centralizer calculation
gives $K=1$.  This proves the assertions.
\end{proof}

\begin{lemma}\label{ext:soluble-interval}
A nontrivial fenced interval cannot be an interval in a finite soluble
group.  Consequently, if
$\operatorname{Inn}(T)\le B\le H\le\operatorname{Aut}(T)$, the interval
$[B,H]$ is not nontrivially fenced.
\end{lemma}
\begin{proof}
Lemma~\ref{frm:labels} would give a nonabelian simple composition factor
of the soluble top group.  For the second assertion, quotient by
$\operatorname{Inn}(T)$, which is contained in the bottom $B$.
The resulting interval lies in the soluble group
$\operatorname{Out}(T)$, by Schreier's Theorem.
\end{proof}

\subsection{Minimality in both orientations}

Let $\mathcal L$ be the lattice in
Proposition~\ref{lat:decorated-lattice}.  Suppose that one of $\mathcal L$
and $\mathcal L^{\mathrm{op}}$ is a finite group interval.  Among all
representations of either orientation choose
\[
             [D,G]\cong\mathcal M,
       \qquad \mathcal M\in\{\mathcal L,\mathcal L^{\mathrm{op}}\},
\]
with $|G|$ minimal.  Minimizing over both orientations is the same
kind of reduction used in the subgroup-interval and signalizer setting
of \cite[Theorem~3 and Section~6]{Aschbacher2008}; see also the comparison in
\cite[Section~4]{PalfyTwisted}.  We give the needed complementary-socle
argument in full.  Quotienting by $\operatorname{core}_G(D)$ preserves
the interval; hence this core is trivial.  The global fence and
Lemma~\ref{frm:labels} give
\begin{equation}\label{ext:top-socle}
 S=\operatorname{soc}(G)=\prod_{i\in I}T_i\cong T^m,
 \qquad G=DS,
\end{equation}
where $T$ is nonabelian simple and $D$ acts transitively on the factor
set $I$.  Proposition~\ref{lat:decorated-lattice} applies to the chosen
orientation $\mathcal M$.

\begin{proposition}\label{ext:trivial-intersection}
In this minimal representation, $D\cap S=1$.
\end{proposition}
\begin{proof}
Put $E=D\cap S$.  For every $Y\in[D,G]$, the identity $G=DS$ gives
\[
                  Y=D(Y\cap S).
\]
Thus $Y\mapsto Y\cap S$ is an isomorphism from $[D,G]$ to the interval
of $D$-invariant subgroups between $E$ and $S$; its inverse is
$K\mapsto DK$.

First suppose that $E$ is subdirect in $S$.  By the subdirect-power
description in Lemma~\ref{frm:subdirect}, its simple factors are
full diagonal strips on a partition $\Pi$ of $I$.  Every subgroup
between $E$ and $S$ is subdirect.  Its strips split the blocks of $\Pi$,
and their diagonal identifications are forced by their containment of
$E$.  Therefore the preceding interval is anti-isomorphic to the
lattice of $D$-invariant equivalence relations refining $\Pi$.

Fix $i_0\in I$, let $A=D_{i_0}$, and let $B$ be the stabilizer in $D$
of the block of $\Pi$ containing $i_0$.  The block correspondence for a
transitive action identifies these equivalence relations with $[A,B]$:
the block containing $i_0$ associated with $U\in[A,B]$ is $U i_0$,
and all its other blocks are its $D$-translates.  Conversely, the
setwise stabilizer of the block containing $i_0$ recovers $U$.
It follows that $\mathcal M^{\mathrm{op}}\cong[A,B]$.  But
$|B|\le |D|<|G|$, contrary to the choice of $G$ over both orientations.

The coordinate projections of $E$ have equal order, since $D$ is
transitive on $I$.  If they are proper and nontrivial,
Proposition~\ref{frm:strict-product}, applied with the possibly unlabelled
lower endpoint $D$, and the global fence imply that $E$ is the product
of those projections.  Choose a proper original Boolean vertex $Y$
of the decorated lattice.  Its forward label equals the top label,
so $Y\cap S$ is subdirect by Proposition~\ref{frm:equal-transport}.
It is proper, because $Y<G$.

A product of nontrivial coordinate subgroups cannot lie in a proper
subdirect subgroup: the latter has a diagonal link between two
coordinates, whereas the former contains an element that is nontrivial
in one of those coordinates and trivial in the other.  This contradicts
$E\le Y\cap S$.  The only remaining possibility is that all coordinate
projections of $E$ are trivial, and then $E=1$.
\end{proof}

\subsection{The extension dictionary}

We give the parametrization with its group actions and order relations
explicit.  It is the induced-coordinate form of
\cite[Proposition~2.2 and Lemma~2.3]{PalfyTwisted}.

\begin{proposition}\label{ext:dictionary}
Let a finite group $D$ act on $S=\prod_{i\in I}T_i$, where all $T_i$
are isomorphic to a nonabelian simple group $T$, and suppose that its
action on $I$ is transitive.  Fix $i_0\in I$, set $A=D_{i_0}$, identify
$T_{i_0}$ with $T$, and let
$\alpha:A\to\operatorname{Aut}(T)$ be the induced action.

The $D$-invariant subdirect subgroups of $S$ are in bijection with pairs
$(U,\beta)$, where
\[
 A\le U\le D,\qquad
 \beta:U\longrightarrow\operatorname{Aut}(T),\qquad
 \beta|_A=\alpha.
\]
Write $S(U,\beta)$ for the corresponding subgroup.  Then
\begin{equation}\label{ext:dictionary-order}
 S(U,\beta)\le S(V,\gamma)
 \quad\Longleftrightarrow\quad
 V\le U\ \hbox{ and }\ \beta|_V=\gamma.
\end{equation}
In particular, the interval of $D$-invariant subgroups from
$S(U,\beta)$ to $S$ is anti-isomorphic to $[A,U]$, by restricting
$\beta$.
\end{proposition}
\begin{proof}
Use left conjugation for the given action and let $\pi_0:S\to T$ be
the chosen coordinate projection.  The map
\[
 s\longmapsto f_s,\qquad f_s(d)=\pi_0(d\cdot s),
\]
identifies $S$ with
\[
 \mathcal F_\alpha
 =\{f:D\to T:f(ad)=\alpha(a)(f(d))
                       \text{ for }a\in A,d\in D\}.
\]
Indeed, values on representatives of $A\backslash D$ are independent
coordinates of $S$.  Under this identification the action is
$(g\cdot f)(d)=f(dg)$.

An extension $\beta$ defines
\[
 \mathcal F_\beta
 =\{f:D\to T:f(ud)=\beta(u)(f(d))
                       \text{ for }u\in U,d\in D\}.
\]
This is a $D$-invariant subgroup of $\mathcal F_\alpha$, with independent
values on the cosets in $U\backslash D$.  Consequently
$\mathcal F_\beta\cong T^{[D:U]}$, and it projects onto every coordinate
of $\mathcal F_\alpha$.  Set $S(U,\beta)=\mathcal F_\beta$ in these
coordinates.

Conversely, let $K\le\mathcal F_\alpha$ be a $D$-invariant subdirect
subgroup.  Its diagonal-strip partition is invariant under the
transitive action on $A\backslash D$.  The block containing $A$ has
stabilizer $U$ for some $A\le U\le D$.  Projection at $1$ identifies
the diagonal factor on that block with $T$.  The induced action of $U$
on this factor gives a homomorphism
$\beta:U\to\operatorname{Aut}(T)$ extending $\alpha$.
For $f\in K$ and $u\in U$ one has $f(u)=\beta(u)(f(1))$.
Applying this identity to $d\cdot f$ gives
\[
                     f(ud)=\beta(u)(f(d)).
\]
Thus $K\le\mathcal F_\beta$.  Both groups have one freely varying
diagonal factor for each coset of $U$, so they are equal.

Containment of two subdirect groups reverses refinement of their strip
partitions.  Hence $S(U,\beta)\le S(V,\gamma)$ forces $V\le U$.
For $v\in V$, varying $f(1)$ over $T$ in the two defining equations
then gives $\beta(v)=\gamma(v)$.  The converse follows immediately
from those equations.  Finally, every subgroup containing a subdirect
subgroup is subdirect, so all members of the stated interval have
already been included.
\end{proof}

Apply Proposition~\ref{ext:dictionary} to
\eqref{ext:top-socle}.  By Proposition~\ref{ext:trivial-intersection},
the full interval $[D,G]$ is the lattice of $D$-invariant subgroups of
$S$, via $Y\mapsto Y\cap S$.  Whenever $Y\cap S$ is subdirect, write
$(U_Y,\beta_Y)$ for its extension.  In particular,
\begin{equation}\label{ext:actual-filter}
                  [Y,G]^{\mathrm{op}}\cong[A,U_Y].
\end{equation}
For nested filters these are compatible actual group intervals: the
larger domain's homomorphism restricts to the smaller domain's
homomorphism.  No representation of the entire dual lattice is asserted
by \eqref{ext:actual-filter}.

\subsection{The detector and the common image}

\begin{proposition}\label{ext:detector}
In the minimal representation fixed above, the induced factor-stabilizer
action $\alpha$ satisfies
$\alpha(A)\ge\operatorname{Inn}(T)$.
\end{proposition}
\begin{proof}
Let $X$ be the detector vertex of
Proposition~\ref{lat:decorated-lattice}.  Its forward type equals $T$, so
$E_X=X\cap S$ is subdirect in $S$.  The action of $D$ on the simple
factors of $E_X$ is transitive.  Its factor-stabilizer action is
$\beta_X:U_X\to\operatorname{Aut}(T)$ by the construction in
Proposition~\ref{ext:dictionary}.

Suppose that this action preserved a subgroup $1<P<T$.  Transporting
$P$ among the factors of $E_X$ would give a $D$-invariant product $Q$
of proper nontrivial coordinate subgroups of $E_X$.  Thus $DQ$ is
strictly between $D$ and $X$.  Choose a coatom $Y$ of $[D,X]$ containing
$DQ$.  The detector's coatom property gives a designated vertex
$d\le Y$ whose intersection $d\cap S$ is subdirect in the global
coordinates.  Hence $Y\cap S$ projects onto every simple factor of
$E_X$: projection onto such a factor can be identified with projection
onto any one global coordinate in its diagonal support.
Thus $Y\cap S$ is a proper subdirect subgroup of $E_X$ containing the
product $Q$, which is impossible by the diagonal-link argument in
Proposition~\ref{ext:trivial-intersection}.

The automorphism lemma recalled by P\'alfy
\cite[Lemma~2.11]{PalfyTwisted} states that a subgroup of
$\operatorname{Aut}(T)$ not containing $\operatorname{Inn}(T)$
preserves some proper nontrivial subgroup of $T$.
It follows that
\[
                 \beta_X(U_X)\ge\operatorname{Inn}(T).
\]

The reverse detector interval is the actual interval $[A,U_X]$ in
\eqref{ext:actual-filter}, and it is fenced.  Set $K=\ker\beta_X$.
Lemma~\ref{frm:normal-test} gives either $AK=U_X$ or $K\le A$.
The first possibility gives
$\beta_X(U_X)=\beta_X(A)=\alpha(A)$, as required.

In the second possibility $K\le C_X=\operatorname{core}_{U_X}(A)$,
and $[A,U_X]$ is represented as
\[
             [\alpha(A),\beta_X(U_X)].
\]
Write $H=\beta_X(U_X)$ and $B=\alpha(A)$.  If
$\operatorname{core}_H(B)$ is nontrivial, it contains
$\operatorname{Inn}(T)$ by Lemma~\ref{ext:almost-simple}.  Quotienting
by that core leaves a soluble top, contradicting
Lemma~\ref{ext:soluble-interval}.
Otherwise $C_X=K$, and $U_X/C_X\cong H$ has socle $T$ with multiplicity
one.  The reverse detector supplies a strictly smaller labelled domain
in $[A,U_X]$ whose simple type is equal to that of $U_X/C_X$.
The equality is supplied by the reverse tests in
Proposition~\ref{lat:decorated-lattice}, applied through
\eqref{ext:actual-filter}.  Proposition~\ref{frm:equal-transport} would make
its positive socle multiplicity strictly smaller than one.  This final
contradiction excludes $K\le A$.
\end{proof}

\begin{lemma}\label{ext:constant-image}
With $\alpha$ as in Proposition~\ref{ext:detector}, let $A<U\le D$ and
let $\beta:U\to\operatorname{Aut}(T)$ extend $\alpha$.
Suppose that $[A,U]$ is fenced.  Then
\[
                  U=A\ker\beta,
       \qquad \beta(U)=\alpha(A).
\]
\end{lemma}
\begin{proof}
The normal-subgroup test gives the asserted supplementation unless
$\ker\beta\le A$.  In the latter case $[A,U]$ is isomorphic to
$[\alpha(A),\beta(U)]$.  Proposition~\ref{ext:detector} puts the inner
group in its bottom, contrary to Lemma~\ref{ext:soluble-interval}.
Supplementation gives the equality of images.
\end{proof}

\begin{proposition}[The Boolean extension configuration]
\label{ext:reduction}
If either orientation of the decorated lattice has a finite-group
interval representation, there are finite groups $A\le D$, finite
nonabelian simple groups $T,R$, a homomorphism
$\alpha:A\to\operatorname{Aut}(T)$, and extensions
\[
 \beta_J:U_J\to\operatorname{Aut}(T)
       \qquad(\varnothing\ne J\subsetneq\{1,\ldots,4\})
\]
with the following properties.  Set $U_\varnothing=A$ and
$\beta_\varnothing=\alpha$.
\begin{enumerate}[label=\textup{(E\arabic*)}]
\item\label{ext:e1}
For nonempty proper $J$, $[A,U_J]$ is fenced.  If
$C_J=\operatorname{core}_{U_J}(A)$, the group $U_J/C_J$ has unique
minimal normal subgroup
\[
 N_J/C_J\cong R^{I_J},\qquad U_J=AN_J.
\]
\item\label{ext:e2}
If $J\subsetneq K\subsetneq\{1,\ldots,4\}$, then $U_J<U_K$ and
$\beta_K|_{U_J}=\beta_J$.  For nonempty $J$, the equal-type transport
of Proposition~\ref{frm:equal-transport} applies: in particular
\[
 \frac{U_J\cap N_K}{C_K}\ \cong\ \frac{N_J}{C_J},
 \qquad N_J=(U_J\cap N_K)C_J,
\]
and the canonical surjections $I_K\to I_J$ are $U_J$-equivariant and compose.
Whenever a nonempty family of proper subsets has proper union $K$, its domains
generate $U_K$.
\item\label{ext:e3}
All images $\beta_J(U_J)$ equal
$H=\alpha(A)\ge\operatorname{Inn}(T)$.
\item\label{ext:e4}
If a family of proper nonempty subsets has union $\{1,\ldots,4\}$,
its homomorphisms have no common extension to the subgroup of $D$
generated by their domains.
\end{enumerate}
\end{proposition}
\begin{proof}
Use the minimal representation already chosen.  Index the original
Boolean component in its reverse order by subsets $J$ of
$\{1,\ldots,4\}$.  Let $Y_J$ be the corresponding forward group, so
$Y_\varnothing=G$ and $Y_{\{1,\ldots,4\}}=D$.
Every proper nontrivial $Y_J$ has the top forward type and hence a
subdirect intersection with $S$.  Proposition~\ref{ext:dictionary}
provides $(U_J,\beta_J)$.

The reverse labels and tests in
Proposition~\ref{lat:decorated-lattice} apply in the actual filters
\eqref{ext:actual-filter}.  They give the fences and equal simple types
on comparable proper Boolean nodes.  All the nonempty proper nodes
therefore have one common simple type $R$: connect two atoms through
their proper two-atom join, and connect any other node to an atom below
it.  Lemma~\ref{frm:labels} and Proposition~\ref{frm:equal-transport} prove
\ref{ext:e1} and the transport assertions in \ref{ext:e2}.
The order relation and restrictions follow from
\eqref{ext:dictionary-order}.

For the assertion about a proper join, let $K$ be the union of a
nonempty family of index sets, and put
$V=\langle U_J:J\text{ in the family}\rangle$.
Then $A\le V\le U_K$, and $\beta_K|_V$ is an extension of all the
$\beta_J$.  Its subdirect group is contained in every $Y_J\cap S$
and contains $Y_K\cap S$.  Original Boolean meets are preserved in the
decorated lattice, so the intersection of those $Y_J$ is exactly
$Y_K$.  The subdirect groups are equal; the dictionary gives $V=U_K$.

Proposition~\ref{ext:detector} and Lemma~\ref{ext:constant-image} give
\ref{ext:e3}.  Finally, a common extension for a family with full union
would give a nontrivial subdirect subgroup of $S$ contained in all the
$Y_J\cap S$.  Their intersection is $D\cap S=1$, because their reverse
Boolean join is the top.  This proves \ref{ext:e4}.
\end{proof}
